\documentclass[12pt]{article}
\usepackage[T1]{fontenc}
\usepackage[utf8]{inputenc}
\usepackage{lmodern}
\usepackage{textcomp}
\usepackage{enumitem}
\usepackage{geometry}
\usepackage{graphicx}
\usepackage{amsmath, amssymb}
\geometry{margin=1in}
\begin{document}
\begin{center}
Constitutional Law - Graduate Level Assessment 15 \
Total Marks: 40 \
Answer all questions.
\end{center}

\section*{Multiple Choice (10 marks)}
\begin{enumerate}[label=\arabic*.]
\item Which of the following propositions would be rejected by those who describe themselves as critical legal theorists?
\begin{enumerate}[label=(\alph*)]
    \item Law is unstable.
    \item Law is determinate.
    \item Law reflects economic power.
    \item Law is politics.
    \item Law is a social construct.
\end{enumerate}
\item Why does Parfit oppose equality?
\begin{enumerate}[label=(\alph*)]
    \item He claims that by giving priority to the needs of the poor, we can increase equality.
    \item He rejects the idea of equality altogether.
    \item He believes that equality is not a realistic goal.
    \item He argues than an unequal society is inevitable.
    \item He suggests that societal structures inherently support inequality.
\end{enumerate}
\item The \_\_\_\_\_\_\_\_ School of jurisprudence postulates that the law is based on what is "correct."
\begin{enumerate}[label=(\alph*)]
    \item Legal Pragmatism
    \item Legal Formalism
    \item Comparative
    \item Analytical
    \item Sociological
\end{enumerate}
\item Why is it important to separate the concept of punishment from its justification?
\begin{enumerate}[label=(\alph*)]
    \item Because the concept of punishment has evolved over time.
    \item Because punishment can be justified in multiple ways.
    \item Because any definition of punishment should be value-neutral.
    \item Because the practice of punishment is separate from its justification.
    \item Because the justification of punishment varies across cultures.
\end{enumerate}
\item In which of the following stages does an indigent person not have the Sixth Amendment right to counsel?
\begin{enumerate}[label=(\alph*)]
    \item Execution of the arrest warrant
    \item investigative surveillance
    \item Police interrogation prior to arrest
    \item Post-charge lineups
    \item Preliminary Hearing
\end{enumerate}
\end{enumerate}

\section*{True / False (10 marks)}
\textit{Answer True or False}
\begin{enumerate}[label=\arabic*.]
\setcounter{enumi}{5}
\item True or False: Hume argued that natural law effectively secures the state's stability against external threats.
\item True or False: The protective principle of jurisdiction applies solely to crimes committed within the territory of a state.
\item True or False: The international recognition of human rights after World War II significantly contributed to the revival of natural law in the 20 th century.
\item True or False: Partial performance of a contract does not terminate the contract by operation of law.
\item True or False: A three-fourths vote in the Senate is required to convict and remove a president from office after impeachment by the House.
\end{enumerate}

\section*{Long Form Response (20 marks)}
\textit{Answer each question with a clear and complete explanation.}
\begin{enumerate}[label=\arabic*.]
\setcounter{enumi}{10}
\item Discuss the admissibility of evidence in court by analyzing why a diary entry from the insured's mother, regarding the insured's birth date, may be deemed inadmissible in a trial concerning life insurance proceeds.
\item Discuss the relationship between legal positivism and morality, specifically addressing the claim that it regards morals and law as inseparable. Analyze its implications for constitutional law.
\end{enumerate}

\vfill
\noindent\textit{End of Paper}
\end{document}